\section{Affine polarization and the complete square sector}
\label{ag:sec:polarization}

\subsection{A constant-coefficient quadratic symbol}
Let $p,q,t$ be formal parameters and make the differential substitution
\begin{equation}\label{ag:eq:affine}
 u_0=q+px+\varepsilon v(x),\qquad
 u_1=p+\varepsilon v'(x),\qquad
 u_j=\varepsilon v^{(j)}(x)\quad(j\ge2).
\end{equation}
Put $W=e^{t(q+px)}$ and $\kappa=tp$, so $W'=\kappa W$.
For a quadratic differential polynomial in $v$ with constant
coefficients, polarization introduces two derivative variables
$\xi_1,\xi_2$. Modulo a total derivative with weight $W$, the relation is
\[
 \xi_1+\xi_2+\kappa=0.
\]
Set $\xi_1=\xi$, $\xi_2=-\xi-\kappa$, and
\begin{equation}\label{ag:eq:eta}
 \eta=-\xi(\xi+\kappa).
\end{equation}
The symbol of $(v^{(r)})^2$ is $\eta^r$, while the symbol of $vv'$ is
$-\kappa/2$. A weighted derivative is annihilated because its symbol is
multiplied by $\xi_1+\xi_2+\kappa$.

\begin{lemma}[Uniqueness of the square symbol]\label{ag:lem:symbol}
If a finite sum $W\sum_{r=0}^M c_r(v^{(r)})^2$ is a weighted total
derivative of a constant-coefficient quadratic differential expression,
then every $c_r$ is zero. The same conclusion holds coefficientwise for
formal series in $t$.
\end{lemma}
\begin{proof}
Its symbol is $\sum c_r\eta^r=0$. Since $\eta$ is a polynomial of
degree two in $\xi$, distinct powers have different degrees and are
linearly independent. The formal statement follows at every finite
coefficient of $t$. Conversely, the symmetric symbols invariant under
$\xi\mapsto-\xi-\kappa$ form the polynomial ring in $\eta$, so this
symbol also supplies the usual diagonal reduction of quadratic forms.
\end{proof}

Lemma~\ref{ag:lem:appell-primitive} ensures that this is exactly the
quotient needed after \eqref{ag:eq:affine}: the quadratic part of the
primitive of $\mathcal C-\NF\mathcal C$ is $W$ times an expression
whose coefficients do not depend on $x$. An arbitrary $x$-dependent
antiderivative, which would invalidate the symbol argument, is not
being allowed.

\subsection{The symbol of the centered Gamma exponential}
Define the even differential series
\[
 K_t(z)=\frac{2\sinh(tz/2)}{z},\qquad K_t(0)=t.
\]
The exponent in \eqref{ag:eq:centered} is $K_t(\D)u_0$. Its affine
part is $t(q+px)$, so the quadratic perturbation of $\mathcal C$ is
\begin{equation}\label{ag:eq:quadratic-density}
 [\varepsilon^2]\mathcal C=\frac{W}{2}\bigl(K_t(\D)v\bigr)^2.
\end{equation}
Now set $y=t^2p/2$ and $h=t(\xi+\kappa/2)$; then
$h^2=y^2-t^2\eta$. The product-to-sum identity for hyperbolic sines
gives the exact symbol
\begin{align}
 \frac12K_t(-\xi-\kappa)K_t(\xi)
 &=\frac{\cosh y-\cosh\sqrt{y^2-t^2\eta}}{\eta}\notag\\
 &=\sum_{r\ge0}\frac{(-1)^rt^{2r+2}}{(r+1)!}
           \piy^{\,r+1}\cosh y\,\eta^r.
 \label{ag:eq:master-symbol}
\end{align}
The square root introduces no branch: $\cosh\sqrt z$ means its entire
series $\sum_{n\ge0}z^n/(2n)!$. The numerator is divisible by $\eta$.

On the other hand, normal support and the Appell factor imply
\eqref{ag:eq:normal-sector}, initially with unknown $f_0,g_r$.
No normal monomial has exactly one factor of index at least two.
If it has exactly two, their indices must coincide. All other such
monomials contribute only at order $\varepsilon^3$ or higher.
The quadratic perturbation of its low-index part is
\[
 W\left\{\frac{t^2f_0}{2}v^2+\frac{t^3f_0'}{2}vv'
                         +\frac{t^4f_0''}{8}(v')^2\right\}.
\]
After replacing the mixed term by its symbol, the coefficients of
$\eta^0$ and $\eta^1$ are
\[
 \frac{t^2}{2}(f_0-yf_0'),\qquad \frac{t^4}{8}f_0''.
\]
For $r\ge2$ the coefficient of $\eta^r$ is simply
$t^{2r+2}g_r(y)$. Lemma~\ref{ag:lem:symbol} and
\eqref{ag:eq:master-symbol} therefore prove the derivative formula
for $g_r$ in Theorem~\ref{ag:thm:main}.

The coefficient of $\eta^0$ gives
\begin{equation}\label{ag:eq:f0-ode}
 f_0-yf_0'=\frac{\sinh y}{y}.
\end{equation}
The normal form excludes a term linear in $u_1$, hence a term linear
in $y$. With $f_0(0)=1$, the unique solution is
\eqref{ag:eq:f0-intro}. The $\eta^1$ equation agrees with its second
derivative and supplies an independent normalization check.

\subsection{The beta representation and explicit coefficients}
\begin{proposition}[Entire kernels and their recurrences]\label{ag:prop:kernel}
For every $r\ge0$ define
$g_r(y)=(-1)^r\piy^{r+1}\cosh y/(r+1)!$. Then
\begin{equation}\label{ag:eq:beta-kernel}
 g_r(y)=c_r\int_0^1(1-v^2)^r\cosh(yv)\dd v,
 \qquad c_r=\frac{(-1)^r}{2^{2r+1}r!(r+1)!}.
\end{equation}
They are even entire functions and satisfy
\begin{align}
 \sum_{r\ge0}g_r(y)w^r
   &=\frac{\cosh y-\cosh\sqrt{y^2-w}}{w},
       \label{ag:eq:bivariate}\\
 g_{r+1}(y)&=-\frac1{r+2}\piy g_r(y),
       \label{ag:eq:kernel-rec}\\
 yg_r''+(2r+2)g_r'-yg_r&=0,
 \qquad g_r(0)=\frac{(-1)^r}{(2r+2)!}.
       \label{ag:eq:kernel-ode}
\end{align}
In \eqref{ag:eq:bivariate}, the terms $r=0,1$ are symbol coefficients;
they are not $f_0$ and a nonexistent normal $u_1$ term.
\end{proposition}
\begin{proof}
Expanding the hyperbolic cosine in \eqref{ag:eq:beta-kernel} is uniformly
convergent on compact sets. For every even $k\ge0$,
\[
 \int_0^1v^k(1-v^2)^r\dd v
 =\frac12 B\left(\frac{k+1}{2},r+1\right)
 =\frac{2^r r!}{\prod_{j=0}^r(k+2j+1)}.
\]
Thus
\begin{equation}\label{ag:eq:kernel-coeff}
 k![y^k]g_r(y)=
 \frac{(-1)^r}{2^{r+1}(r+1)!\prod_{j=0}^r(k+2j+1)}
 \quad(k\text{ even}),
\end{equation}
and odd coefficients vanish. Comparing with the differentiated cosine
series proves the beta identity and the ${}_0F_1$ expression.
Taylor expansion in $w$ proves \eqref{ag:eq:bivariate}; differentiation
proves \eqref{ag:eq:kernel-rec}. Substituting the coefficients proves
\eqref{ag:eq:kernel-ode}. Finally, expand $e^{tu_0}$ and set
$y=u_1t^2/2$ to obtain \eqref{ag:eq:coefficient}.
\end{proof}

The first two high-index kernels are elementary:
\begin{align}
 g_2(y)&=\frac{(y^2+3)\sinh y-3y\cosh y}{48y^5},
       \label{ag:eq:g2}\\
 g_3(y)&=-\frac{(y^3+15y)\cosh y-(6y^2+15)\sinh y}{384y^7}.
       \label{ag:eq:g3}
\end{align}
At the origin these formulas should be evaluated by their even Taylor
series or by ${}_0F_1$, rather than by subtracting nearly equal floating
numbers. The recurrence generates every higher elementary hyperbolic
expression, while the beta form is preferable for the analytic bridge.
